\section{历年真题全流程精讲}
本章把前面的构件放回真实题目。每题先压缩规则，再明确边界、选择结构、安排实现顺序，最后给出完整C++17程序。题目入口保留完整原题；下文用自己的规则表和小例子讲解。可先合上代码，按函数职责写一版，再对照实现检查差异。

\subsection{元素选择器：从缩进文本到祖先链}
\ojpractice{CSP201809C}{元素选择器：第14次认证，2018年9月，第3题}
\subsubsection{题面瘦身}
\begin{enumerate}
\item 文档每行一个元素，行号从1开始。开头每两个点号表示一层缩进；点号后是标签，可附带以井号开头的ID。
\item 标签按忽略大小写的方式比较，ID按原大小写比较。文档结构合法，ID唯一。
\item 单个选择符按标签或ID匹配。多个选择符以空格连接时，最后一个匹配候选节点，之前的依次匹配它的严格祖先，层与层之间可以跨过节点。
\item 每条查询输出匹配数，随后按行号递增输出匹配的节点。一个节点只有符合完整选择器时才被选中。
\end{enumerate}
输入第一行$n,q$，接下来$n$行文档，再接$q$行查询。$n\le100$，$q\le10$，每行至多80字符。下面是用于逐步推演的自拟小例子：
\begin{verbatim}
输入
6 4
HTML
..DIV #Main
....p #A
....DIV
......P #a
..p
DIV p
#Main #a
div div p
#A
输出
2 3 5
1 5
1 5
1 3
\end{verbatim}

\subsubsection{读题避坑}
先数行首点号，再解析后面的字段。整行转小写会把\texttt{\#A}与\texttt{\#a}混在一起，应只把标签及标签选择符转成小写。标签与ID分开存，缺少ID时保留空串。

后代允许跨层。例如\texttt{\#Main \#a}从第5行向上经过第4行后，在第2行找到Main。\texttt{div div p}则要求两个不同的div祖先：第5行有第4、2行，第3行只有第2行。节点匹配最后一个选择符之后，从它的父亲开始继续查找，保证各层使用不同节点。

\subsubsection{数据结构选型}
结构体保存tag、id、parent三个字段。祖先栈按深度保存当前路径：读到深度$d$时，把栈长度截为$d$，栈顶就是父亲，然后压入新节点。例如读完第5行后栈保存第1、2、4、5行；读第6行时深度回到1，截断后只留下第1行，因此父亲恢复为第1行。

匹配时从右向左处理查询。先匹配候选节点，剩下的选择符沿父链向上寻找，遇到匹配的祖先就消费一个选择符。优先选最近祖先，会给剩余选择符留下最多的上方节点：若某个解使用更远的同类祖先，把它换成最近祖先后，更高层的匹配仍然可用。按节点行号从小到大枚举，答案自然有序。

\subsubsection{从子任务到完整解法}
先完成单标签查询，再增加ID查询，随后增加两级与多级后代。原题前六个测试点围绕标签、ID和不同树高，后四个加入后代匹配。父编号建好以后，多级匹配只是让选择符下标多走几次，基础结构可以保留。

设树高$H$、查询选择符数$K$、单次字符串比较最长$C$，查询时间为$O(qn(H+K)C)$，空间为文档与查询文本的大小。按原题范围直接沿祖先链即可完成全部查询。代码中lower\_tag只负责标签转换，matches只负责一个节点，selected负责完整祖先链；这样可以分别检查输入、单点匹配和多级匹配。

\subsubsection{完整参考代码}
\Code{69_selector_real.cpp}

\subsection{化学方程式：递归返回元素计数表}
\ojpractice{CSP201912C}{化学方程式：第18次认证，2019年12月，第3题}
\subsubsection{题面瘦身}
\begin{enumerate}
\item 每行一个合法方程式，等号两边分别由若干化学式用加号连接。
\item 每个化学式前可带正整数系数；元素或括号后可带正整数数量，省略时均为1。
\item 元素名由一个大写字母，或一个大写字母紧跟一个小写字母组成；括号可嵌套。
\item 分别统计等号两边每种元素的原子数，完全相同输出Y，否则输出N。
\end{enumerate}
第一行方程式数量$n\le100$。每行长度不超过1000，无空白，系数无前导零，每边每种元素总数至多$10^9$。任务是检查给定系数是否配平。下列自拟样例覆盖普通系数、括号、嵌套和元素名的区分：
\begin{verbatim}
输入
4
H2+O2=H2O
2Al(OH)3+3H2SO4=Al2(SO4)3+6H2O
K4(ON(SO3)2)2=K4O14N2S4
Co=CO
输出
N
Y
Y
N
\end{verbatim}

\subsubsection{读题避坑}
数字的位置决定作用范围：化学式前的数字乘整个化学式，元素后的数字乘单个元素，右括号后的数字乘括号内全部元素。\texttt{3Ca(OH)2}先得到Ca为1、O为2、H为2，再整体乘3，结果为Ca为3、O和H各为6。

\texttt{Co}是一个元素，\texttt{CO}是C与O两个元素。读一个大写字母后，只再接收紧随的小写字母。多位数字连续扫描，遇到分隔符时把指针留给调用者。原题保证语法合法，可以按已确定的语法前进。

\subsubsection{数据结构选型}
用\texttt{map<string, long long>}保存“元素名到数量”的映射。把表达式求值的返回值从整数换成计数表：元素返回单项表，括号返回内层计数表，连接操作把同名元素数量相加，数量与系数对整张表逐项乘法。

解析分为三层。number读取数字，缺省返回1；formula读取连续的元素或括号项，在右括号、加号、等号或行尾处停止；side循环读取“前置系数+formula”，遇到等号或行尾停止。formula遇到左括号后递归，递归返回时消费右括号，再读取紧随其后的数量。共享pos始终指向下一个未处理字符。

以\texttt{Al2(SO4)3}为例：读取Al后数量2，表为Al:2；进入括号得到S:1、O:4；括号后数量3，合并后成为Al:2、S:3、O:12。两个side结果比较时，map会同时比较元素名和数量，键的先后输入顺序不影响结果。

\subsubsection{从子任务到完整解法}
原题前六个测试点逐步加入小写字母、加号和数字；第7、8点加入一层括号；第9、10点允许嵌套。先写一个扫描器处理元素名、数量与化学式前系数，再把括号部分交给同一个formula递归，便能沿用计数合并操作。

设行长$L$、最多$E$种元素、括号深度$D$。每次合并至多处理$E$项，map更新耗时$O(\log(E+1))$，用$O(LE\log(E+1))$作上界；活动递归表占$O(DE)$空间。元素名长度至多2，可能的键至多$26\times27$种。按题目计数上界使用long long保存中间乘法与总量。

\subsubsection{完整参考代码}
\Code{70_chem.cpp}

\subsection{角色授权：把多表关联变成集合运算}
\ojpractice{CSP202206C}{角色授权：第26次认证，2022年6月，第3题}
\subsubsection{题面瘦身}
\begin{enumerate}
\item 每个角色包含三张允许清单：操作、资源种类、资源名称。操作和种类中的星号匹配任意值；资源名称清单为空时匹配任意名称。
\item 关联记录把一个角色授予若干用户或用户组。一个角色可以出现在多条关联记录中，同一用户也可以获得多个角色。
\item 每次查询给出用户名、该次所属的用户组、操作、资源种类和资源名称。通过直接用户或任意当前用户组找到候选角色。
\item 只要一个候选角色同时满足三项权限，就输出1；全部候选均不满足时输出0。
\end{enumerate}
输入首先给$n,m,q$。随后$n$个角色，每个按“角色名、操作数及清单、种类数及清单、名称数及清单”读取；接着$m$条关联，格式为“角色名、对象数、若干组类型与名称”，类型u表示用户、g表示组；最后$q$次查询，格式为“用户名、组数及各组名、操作、种类、名称”。

$n,m\le500$，$q\le5000$；每张清单、每条关联对象数与每次查询的组数至多400，字符串长度至多10。下面的小样例中，alice通过组获得读取权限，bob通过用户名获得全部权限：
\begin{verbatim}
输入
2 2 4
reader 1 read 1 file 0
admin 1 * 1 * 0
reader 1 g staff
admin 1 u bob
alice 1 staff read file report
alice 0 read file report
bob 0 delete door front
alice 1 staff write file report
输出
1
0
1
0
\end{verbatim}

\subsubsection{读题避坑}
每次查询使用本次给出的组。样例中alice第一次属于staff，第二次组数为0，第二次就只能使用直接授予alice的角色。用户与用户组使用两张索引表，即使它们名字相同，也各自按u或g解释。

三个权限条件必须由同一个角色满足。假设角色A只允许读文件，角色B只允许写门，用户同时拥有A和B，查询“读门”仍输出0。推导时先在每张权限表中筛出满足条件的角色，再对角色集合取交集，就能保留同一个角色这一要求。

操作与种类用星号表示任意，名称用空清单表示任意。重复出现的关联应合并；查询不存在的用户或组时返回空角色集合。

\subsubsection{数据结构选型}
为$n$个角色分配0到$n-1$的整数编号。用\texttt{bitset<505>}表示一组角色，第$r$位为1表示角色$r$在集合中。例如只有角色0和2时，低三位是101。按位或合并集合，按位与筛出共同角色，any检查集合是否非空。

建立五类倒排索引：用户名到角色集合、组名到角色集合，以及操作、种类、名称各自到允许该值的角色集合。另存三个任意匹配集合。所有索引以字符串为键，以位集合为值。读角色清单时把对应位设为1；读关联时按用户或组把角色位加入对应集合。

对一次查询，先把直接角色和各个当前组的角色取并集，得到$C$；再分别得到满足操作、种类、名称的角色集合$O,T,N$，计算
\[
C\cap O\cap T\cap N.
\]
结果有任意一位为1即授权成功。沿用样例，alice与staff关联得到reader；第二次查询没有组，候选集合变空。bob得到admin，其位同时在三个任意匹配集合中，因此任意行为均通过。

\subsubsection{从70\%子任务到完整解法}
70\%子任务中$n,m\le50$，可以扫描关联记录找角色，再逐个角色扫描三张清单。扩大到满数据后，把清单建成哈希集合可加快单角色查询；进一步使用上述倒排索引和位集合，将多个角色的检查合并成机器字运算。这样每次查询只需读入它的组名，并合并固定长度的角色集合。

设关联对象总数为$A$、角色清单词条总数为$P$、查询组数为$G$，位集合机器字数$W=\lceil505/64\rceil=8$。预处理创建并更新倒排集合；每次查询平均耗时$O((G+1)W)$，另计字符串哈希成本。若索引中不同键共$S$个，位集合存储为$O(SW)$，加上键本身的文本存储。代码中的lookup用find读表，read\_rules集中处理两类通配规则。

\subsubsection{完整参考代码}
\Code{71_rbac.cpp}

\subsection{解压缩：按字节推进的协议解析器}
\ojpractice{CSP202305C}{解压缩：第30次认证，2023年5月，第3题}
\subsubsection{题面瘦身}
\begin{enumerate}
\item 输入$s$个压缩字节，用小写十六进制表示，每行最多8个字节。先把每两个字符转成一个0至255的数。
\item 开头的导引域每字节用低7位记录长度的一部分，最高位为1表示还有下一字节；按低位在前恢复原文长度。
\item 后续逐个解析元素。首字节低两位区分字面量、短回溯和长回溯；字面量追加输入中的原始字节，回溯按偏移和长度追加已有输出中的字节。
\item 回溯允许重叠，刚追加的字节可被继续引用。输出也采用小写十六进制，每行8个字节，末行按实际长度输出。
\end{enumerate}
$s\le2\times10^6$，解压结果长度不超过2MiB，输入保证合法。下面自拟数据表示先写一个字节a，再向前1字节复制8次，得到9个a：
\begin{verbatim}
输入
5
0900611101
输出
6161616161616161
61
\end{verbatim}

\subsubsection{读题避坑}
把“输入字符下标”和“压缩字节下标”区分开：\texttt{61}是一个数值为97的字节。先完成十六进制转换，后续pos每增加1都表示消费一个字节。输出也按字节数每8个换行。

导引域用7位为一组，小端整数用8位为一组。字面量存的是长度减1，短回溯存的是长度减4；写分支前先把公式列成表。令首字节为$b$，低两位为$t$，高六位为$h=\lfloor b/4\rfloor$：
\begin{longtable}{P{2.3cm}P{5.9cm}P{5.8cm}}
\toprule 类型 & 还要读什么 & 还原公式\\\midrule\endhead
$t=0$，$h<60$ & 接下来$l$字节是原文 & $l=h+1$\\
$t=0$，$h\ge60$ & 先读$h-59$字节的小端整数$v$，再读原文 & $l=v+1$\\
$t=1$ & 再读一个字节$x$ & $l=((b\mathbin{\gg}2)\mathbin{\&}7)+4$；$o=(b\mathbin{\gg}5)\times256+x$\\
$t=2$ & 再读两字节小端整数$v$ & $l=h+1$，$o=v$\\
\bottomrule\end{longtable}
原题合法数据包含以上三种低位类型。位运算\texttt{>>}表示右移，\texttt{\&}表示按位与，取低两位可写\texttt{b \& 3}。

最关键的是重叠复制。若已有\texttt{abcde}，偏移2、长度5，就依次从“当前长度减2”的位置取字节，追加d、e、d、e、d。每追加一次，当前长度都会变化。代码每轮重新计算来源，便能覆盖偏移小于长度的情况。

\subsubsection{数据结构选型}
用两个\texttt{vector<unsigned char>}保存压缩字节与已解压字节，pos指向下一压缩字节。little\_endian按位移合并固定数量的字节。导引域读出原文长度后，调用reserve预留输出容量；字面量和回溯都在输出末尾追加。

样例的09是原文长度9；00表示短字面量，长度为1，读到61；11的低两位为01，长度$4+4=8$，紧随01给出偏移1。输出每次从末尾前1字节取61，最终长度恰好为9。整个过程的指针只向前移动。

\subsubsection{从字面量子任务到完整解法}
原题10\%与20\%子任务只含短字面量，40\%加入长字面量，60\%加入短回溯，80\%覆盖两种回溯且原文较短，100\%扩到2MiB。先验证十六进制输入和短字面量，再扩展长度字段，随后加入回溯与大输出。每个阶段都记录pos和当前输出长度，检查每种元素消费了多少输入字节、产生了多少原文字节。

设压缩长度为$s$、原文长度为$n$。每个输入字节读取一次，每个输出字节生成并打印一次，时间$O(s+n)$、空间$O(s+n)$。unsigned int用于位移与字段合并，避免有符号字节扩展；每次取回溯源先复制到一个字节变量，再追加到vector。

\subsubsection{完整参考代码}
\Code{72_decompress.cpp}
